% GENERATED FILE. Edit canonical lessons, then run npm run generate:books.
% canonical-source-hash: fnv1a64:73c35a249e6395cb
% canonical-lessons: GU-C445-malo, GU-C445-agiyo, GU-C445-tamru, GU-C445-tahukvu, GU-C445-cahekvu

\chapter{Nests and Night Sounds}
\label{ch:gu-nests-and-night-sounds}
\hlchaptermodality{445}

\begin{chapteropening}
\textbf{By the end of this chapter:} \emph{I can talk about birds and insects at dawn and at night.}

Complete the last lesson of chapter 445: I can talk about birds and insects at dawn and at night.
\end{chapteropening}

\section[\texorpdfstring{\gu{માળો} (māḷo)}{માળો (māḷo)}]{\gu{માળો} (māḷo) — a nest}

\label{lesson:GU-C445-malo}



\begin{quote}
Before the new one: say the Gujarati for a cluster, then the Gujarati for to be fragrant.
\end{quote}

\subsection*{You'll want to know: \gu{માળો}}
\textbf{\gu{માળો}} — \emph{māḷo} — \textquotedblleft{}a nest\textquotedblright{}.

\begin{quote}
\gu{ચકલીએ} \gu{બારી} \gu{પાસે} \gu{માળો} \gu{બનાવ્યો}.
\end{quote}

\begin{quote}
\emph{caklīe bārī pāse māḷo banāvyo.}
\end{quote}

\begin{quote}
The sparrow built a nest by the window.
\end{quote}

A masculine noun.

\begin{cousinweb}
\textbf{\gu{માળો}} (\emph{māḷo}) goes back to \emph{*māḍa}, \textquotedblleft{}an upper storey\textquotedblright{}, a form CDIAL reconstructs; CDIAL lists a Gujarati word for a loft (CDIAL 9996). How it came to mean a nest is not settled.
\end{cousinweb}

\begin{grammarlens}[title={What you've built}]
One more word for getting around.
\end{grammarlens}

\subsection*{Your turn}
\begin{itemize}
\raggedright
  \item \emph{Say it:} \emph{māḷo}
  \item \emph{Say it:} \emph{māḷo}, once more
  \item \emph{Say it:} \emph{caklīe bārī pāse māḷo banāvyo.} — The sparrow built a nest by the window.
  \item \emph{Recall:} say the Gujarati for a cluster, then the Gujarati for to be fragrant, then say \emph{māḷo} again
\end{itemize}

\begin{tcolorbox}[breakable,colback=teal!4,colframe=teal!35!black,title={Before you move on}]
What does \gu{માળો} mean? (\textquotedblleft{}A nest\textquotedblright{}.) Say it once more.
\end{tcolorbox}
\section[\texorpdfstring{\gu{આગિયો} (āgiyo)}{આગિયો (āgiyo)}]{\gu{આગિયો} (āgiyo) — a firefly}

\label{lesson:GU-C445-agiyo}



\begin{quote}
Before the new one: say the Gujarati for to be fragrant, then the Gujarati for a nest.
\end{quote}

\subsection*{You'll want to know: \gu{આગિયો}}
\textbf{\gu{આગિયો}} — \emph{āgiyo} — \textquotedblleft{}a firefly\textquotedblright{}.

\begin{quote}
\gu{ચોમાસાની} \gu{રાતે} \gu{ખેતરમાં} \gu{આગિયા} \gu{ચમકે} \gu{છે}.
\end{quote}

\begin{quote}
\emph{comāsānī rāte khetarmā̃ āgiyā camke che.}
\end{quote}

\begin{quote}
On monsoon nights fireflies glow in the fields.
\end{quote}

A masculine noun.

\begin{cousinweb}
\textbf{\gu{આગિયો}} (\emph{āgiyo}) is a tadbhava of Sanskrit \emph{agnika}, \textquotedblleft{}fiery\textquotedblright{}; CDIAL lists the Gujarati word, \textquotedblleft{}a glow-worm\textquotedblright{} (CDIAL 57): the little insect with a fire in it.
\end{cousinweb}

\begin{grammarlens}[title={What you've built}]
One more word for getting around.
\end{grammarlens}

\subsection*{Your turn}
\begin{itemize}
\raggedright
  \item \emph{Say it:} \emph{āgiyo}
  \item \emph{Say it:} \emph{āgiyo}, once more
  \item \emph{Say it:} \emph{comāsānī rāte khetarmā̃ āgiyā camke che.} — On monsoon nights fireflies glow in the fields.
  \item \emph{Recall:} say the Gujarati for to be fragrant, then the Gujarati for a nest, then say \emph{āgiyo} again
\end{itemize}

\begin{tcolorbox}[breakable,colback=teal!4,colframe=teal!35!black,title={Before you move on}]
What does \gu{આગિયો} mean? (\textquotedblleft{}A firefly\textquotedblright{}.) Say it once more.
\end{tcolorbox}
\section[\texorpdfstring{\gu{તમરું} (tamrũ)}{તમરું (tamrũ)}]{\gu{તમરું} (tamrũ) — a cricket (the insect)}

\label{lesson:GU-C445-tamru}



\begin{quote}
Before the new one: say the Gujarati for a nest, then the Gujarati for a firefly.
\end{quote}

\subsection*{You'll want to know: \gu{તમરું}}
\textbf{\gu{તમરું}} — \emph{tamrũ} — \textquotedblleft{}a cricket (the insect)\textquotedblright{}.

\begin{quote}
\gu{રાત્રે} \gu{ઘાસમાં} \gu{તમરાં} \gu{બોલે} \gu{છે}.
\end{quote}

\begin{quote}
\emph{rātre ghāsmā̃ tamrā̃ bole che.}
\end{quote}

\begin{quote}
At night crickets chirp in the grass.
\end{quote}

A neuter noun.

\begin{cousinweb}
\textbf{\gu{તમરું}} (\emph{tamrũ}) No dictionary we use records an older form, so its history is not settled.
\end{cousinweb}

\begin{grammarlens}[title={What you've built}]
One more word for getting around.
\end{grammarlens}

\subsection*{Your turn}
\begin{itemize}
\raggedright
  \item \emph{Say it:} \emph{tamrũ}
  \item \emph{Say it:} \emph{tamrũ}, once more
  \item \emph{Say it:} \emph{rātre ghāsmā̃ tamrā̃ bole che.} — At night crickets chirp in the grass.
  \item \emph{Recall:} say the Gujarati for a nest, then the Gujarati for a firefly, then say \emph{tamrũ} again
\end{itemize}

\begin{tcolorbox}[breakable,colback=teal!4,colframe=teal!35!black,title={Before you move on}]
What does \gu{તમરું} mean? (\textquotedblleft{}A cricket (the insect)\textquotedblright{}.) Say it once more.
\end{tcolorbox}
\section[\texorpdfstring{\gu{ટહુકવું} (ṭahukvũ)}{ટહુકવું (ṭahukvũ)}]{\gu{ટહુકવું} (ṭahukvũ) — to call (of a bird)}

\label{lesson:GU-C445-tahukvu}



\begin{quote}
Before the new one: say the Gujarati for a firefly, then the Gujarati for a cricket (the insect).
\end{quote}

\subsection*{You'll want to know: \gu{ટહુકવું}}
\textbf{\gu{ટહુકવું}} — \emph{ṭahukvũ} — \textquotedblleft{}to call (of a bird)\textquotedblright{}.

\begin{quote}
\gu{ચોમાસામાં} \gu{મોર} \gu{ટહુકે} \gu{છે}.
\end{quote}

\begin{quote}
\emph{comāsāmā̃ mor ṭahuke che.}
\end{quote}

\begin{quote}
In the monsoon the peacock calls.
\end{quote}

Past (it …): \gu{તે} \gu{મોર} \gu{ટહુક્યો}, \emph{te mor ṭahukyo}.

\begin{cousinweb}
\textbf{\gu{ટહુકવું}} (\emph{ṭahukvũ}) sounds like the call it names. No dictionary we use records an older form, so its history is not settled.
\end{cousinweb}

\begin{grammarlens}[title={What you've built}]
One more word for getting around.
\end{grammarlens}

\subsection*{Your turn}
\begin{itemize}
\raggedright
  \item \emph{Say it:} \emph{ṭahukvũ}
  \item \emph{Say it:} \emph{ṭahukvũ}, once more
  \item \emph{Say it:} \emph{comāsāmā̃ mor ṭahuke che.} — In the monsoon the peacock calls.
  \item \emph{Recall:} say the Gujarati for a firefly, then the Gujarati for a cricket (the insect), then say \emph{ṭahukvũ} again
\end{itemize}

\begin{tcolorbox}[breakable,colback=teal!4,colframe=teal!35!black,title={Before you move on}]
What does \gu{ટહુકવું} mean? (\textquotedblleft{}To call (of a bird)\textquotedblright{}.) Say it once more.
\end{tcolorbox}
\section[\texorpdfstring{\gu{ચહેકવું} (cahekvũ)}{ચહેકવું (cahekvũ)}]{\gu{ચહેકવું} (cahekvũ) — to chirp; to twitter}

\label{lesson:GU-C445-cahekvu}



\begin{quote}
Before the new one: say the Gujarati for a cricket (the insect), then the Gujarati for to call (of a bird).
\end{quote}

\subsection*{You'll want to know: \gu{ચહેકવું}}
\textbf{\gu{ચહેકવું}} — \emph{cahekvũ} — \textquotedblleft{}to chirp; to twitter\textquotedblright{}.

\begin{quote}
\gu{સવારે} \gu{ચકલીઓ} \gu{ચહેકે} \gu{છે}.
\end{quote}

\begin{quote}
\emph{savāre caklīo caheke che.}
\end{quote}

\begin{quote}
In the morning the sparrows chirp.
\end{quote}

Past (it …): \gu{તે} \gu{ચકલી} \gu{ચહેકી}, \emph{te caklī cahekī}.

\begin{cousinweb}
\textbf{\gu{ચહેકવું}} (\emph{cahekvũ}) No dictionary we use records an older form, so its history is not settled.
\end{cousinweb}

\begin{grammarlens}[title={What you've built}]
One more word for getting around.
\end{grammarlens}

\subsection*{Your turn}
\begin{itemize}
\raggedright
  \item \emph{Say it:} \emph{cahekvũ}
  \item \emph{Say it:} \emph{cahekvũ}, once more
  \item \emph{Say it:} \emph{savāre caklīo caheke che.} — In the morning the sparrows chirp.
  \item \emph{Recall:} say the Gujarati for a cricket (the insect), then the Gujarati for to call (of a bird), then say \emph{cahekvũ} again
\end{itemize}

\begin{tcolorbox}[breakable,colback=teal!4,colframe=teal!35!black,title={Before you move on}]
What does \gu{ચહેકવું} mean? (\textquotedblleft{}To chirp; to twitter\textquotedblright{}.) Say it once more.
\end{tcolorbox}
